\subsection{QUOTIENT SPACES OF METRIZABLE SPACES}

If X is a metrizable space and R is an equivalence relation on X, the quotient space X/R is not necessarily metrizable (even if X is locally compact * and X/R is normal *). However:

PROPOSITION 17. Every Hausdorff quotient space of a compact metrizable space is compact and metrizable.

Equivalently, if \(f\) is a continuous mapping of a compact metrizable space \(\mathbf{X}\) into a Hausdorff space \(\mathbf{X}'\) , then \(f(\mathbf{X})\) is a metrizable subspace of \(\mathbf{X}'\) (Chapter I, § 9, no. 4, Theorem 2, Corollary 4).

Let \(\mathbf{X}\) be a compact metrizable space, and let \(\mathbf{R}\) be an equivalence relation on \(\mathbf{X}\) such that \(\mathbf{X} / \mathbf{R}\) is Hausdorff. Then \(\mathbf{X} / \mathbf{R}\) is compact (Chapter I, § 9, no. 4, Theorem 2), hence by Proposition 16 it is enough to show that the topology of \(\mathbf{X} / \mathbf{R}\) has a countable base. To do this, we use the facts that \(\mathbf{R}\) is closed (Chapter I, § 10, no. 4, Proposition 8) and that the classes mod \(\mathbf{R}\) are compact. Let \(\varphi\) be the canonical mapping of \(\mathbf{X}\) onto \(\mathbf{X} / \mathbf{R}\) , and let \((\mathrm{U}_n)\) be a countable base of the topology of \(\mathbf{X}\) . Let \(z\) be any point of \(\mathbf{X} / \mathbf{R}\) and let \(\mathbf{V}\) be a neighbourhood of \(z\) in \(\mathbf{X} / \mathbf{R}\) ; then \(\overline{\varphi}^1 (\mathbf{V})\) is a neighbourhood in \(\mathbf{X}\) of the compact set \(\overline{\varphi}^1 (z)\) . If \(x\) is any point of \(\overline{\varphi}^1 (z)\) , there is a set \(\mathrm{U}_n\) containing \(x\) and contained in \(\overline{\varphi}^1 (\mathrm{V})\) , and therefore by the Borel-Lebesgue axiom there is a finite open covering \((\mathrm{U}_{n_k})_{1\leqslant k\leqslant r}\) of \(\overline{\varphi}^1 (z)\) such that, if \(\mathbf{W}\) denotes \(\bigcup_{k}\mathrm{U}_{n_k},\mathrm{W}\) is a neighbourhood of \(\overline{\varphi}^1 (z)\) contained in \(\overline{\varphi}^1 (\mathrm{V})\) . Since \(\mathbf{R}\) is closed, it follows that \(\varphi (\mathbf{W})\) is a neighbourhood of \(z\) in \(\mathbf{X} / \mathbf{R}\) , contained in \(\mathbf{V}\) (Chapter I, § 5, no. 4, Proposition 10). Let \(\mathfrak{B}\) denote the set of interiors of sets of the form \(\varphi (\mathbf{W})\) , where \(\mathbf{W}\) runs through the set \(\mathfrak{F}\) of all finite unions of sets of the form \(U_{n}\) ; we have then shown that B is a base of the topology of X/R, and since F is countable, so is B.

PROPOSITION 18. Let X be a complete metric space, let R be an open equivalence relation on X such that X/R is Hausdorff, and let \(\varphi: X \to X/R\) be the canonical mapping. Then if K is any compact subset of X/R, there is a compact subset \(K'\) of X such that \(\varphi(K') = K\) .

Let \(\mathfrak{B}_1\) be the set of all open balls of radius \(1/2\) in X. As B runs through \(\mathfrak{B}_1\) , the sets \(\varphi(B)\) form an open covering of K, and therefore there exists a finite number of points \(x_1, \ldots, x_m\) of X such that the images under \(\varphi\) of the open balls with radius \(1/2\) and centre \(x_i\) ( \(1 \leqslant i \leqslant m\) ) form an open covering of K. Let \(H_1 = \{x_1, \ldots, x_m\}\) and suppose that we have defined a finite set \(H_i\) , for \(1 < i \leqslant n\) , such that:

\begin{enumerate}
\def\labelenumi{(\roman{enumi})}
\item
  \(\mathbf{H}_i\subset \mathbf{H}_{i + 1}\) and each point of \(\mathbf{H}_{i + 1}\) is at a distance \(<  \mathrm{I} / 2^{i}\) from \(\mathbf{H}_i\) , for \(\mathbf{I}\leqslant i\leqslant n - \mathbf{I};\)
\item
  the images under \(\varphi\) of the open balls of radius \(1/2^i\) , with centres at the points of \(\mathbf{H}_i\) , form an open covering of \(\mathbf{K}\) , for \(1 \leqslant i \leqslant n\) .
\end{enumerate}

Let \(\mathfrak{B}_{n+1}\) be the set of all open balls of radius \(1/2^{n+1}\) whose centre \(x\) is such that \(d(x, H_n) < 1/2^n\) ( \(d\) being the metric on X). The properties of \(H_n\) show that the sets \(\varphi(B)\) , for \(B \in \mathfrak{B}_{n+1}\) , form an open covering of K; hence there is a finite set \(L_{n+1} \subset X\) such that the images under \(\varphi\) of the open balls of radius \(1/2^{n+1}\) whose centre belongs to \(L_{n+1}\) form an open covering of K. Taking \(H_{n+1} = H_n \cup L_{n+1}\) , we see that we can define inductively an infinite sequence \((H_n)\) of finite subsets of X with properties (i) and (ii) above. Let \(H = \bigcup_{n} H_n\) , and let us show that H is precompact. For each \(p > 0\) and each point \(z_{n+p} \in H_{n+p}\) , there exists a sequence of points \(z_{n+i} \in H_{n+i}\) ( \(0 \leqslant i \leqslant p - 1\) ) such that

\[
d \left(z _ {n + i}, z _ {n + i + 1}\right) <   \mathrm{I} / 2 ^ {n + i} \quad \text { for } \quad \mathrm{o} \leqslant i \leqslant p - \mathrm{I};
\]

it follows that \(d(z_n, z_{n+p}) \leqslant \sum_{i=0}^{p-1} \mathrm{I}/2^{n+i} \leqslant \mathrm{I}/2^{n-1}\) , and consequently \(d(y, \mathbf{H}_n) \leqslant \mathrm{I}/2^{n-1}\) for all \(y \in \mathbf{H}\) , which proves our assertion. Since \(\mathbf{X}\) is complete, \(\overline{\mathbf{H}}\) is compact, hence \(\varphi(\overline{\mathbf{H}})\) is compact. Next, let us show that \(\mathbf{K} \subset \varphi(\overline{\mathbf{H}})\) . If \(z \in \mathbf{K}\) , then by definition \(d(\overline{\mathbf{H}}_n, \overline{\varphi}^1(z)) \leqslant \mathrm{I}/2^n\) for all \(n\) , and therefore \(d(\overline{\mathbf{H}}, \overline{\varphi}^1(z)) = 0\) ; but \(\overline{\varphi}^1(z)\) is closed and \(\overline{\mathbf{H}}\) is compact, so that this implies \(\overline{\mathbf{H}} \cap \overline{\varphi}^1(z) \neq \emptyset\) (no. 2, Remark following Proposition 3); hence the assertion. Thus if \(\mathbf{K}' = \overline{\mathbf{H}} \cap \overline{\varphi}^1(\mathbf{K})\) , then \(\mathbf{K}'\) is closed in \(\overline{\mathbf{H}}\) and therefore compact, and from what has been proved we have \(\varphi(\mathbf{K}') = \mathbf{K}\) . Q.E.D.

